\begin{center}
\textit{Entgegengesetzte Formen.}
\end{center}

\Article{159}
\textit{Enthält die Form $F$ die Form $F'$ und diese wiederum die Form $F''$, so enthält die Form $F$ auch die Form $F''$.}

Die Unbestimmten der Formen $F,F',F''$ seien der Reihe nach $x,y$, $x',y'$, $x'',y''$, und $F$ gehe in $F'$ über, wenn man
\[
x=\alpha x'+\beta y'.\quad y=\gamma x'+\delta y'
\]
setzt, und $F'$ in $F''$, wenn man
\[
x'=\alpha'x''+\beta'y'',\quad y'=\gamma'x''+\delta'y''
\]
setzt. Dann lässt sich $F$ offenbar in $F''$ überführen durch
\[
x=\alpha(\alpha'x''+\beta'y'')+\beta(\gamma'x''+\delta'y''),\quad
y=\gamma(\alpha'x''+\beta'y'')+\delta(\gamma'x''+\delta'y'')
\]
oder
\[
x=(\alpha\alpha'+\beta\gamma')x''+(\alpha\beta'+\beta\delta')y'',\quad
y=(\gamma\alpha'+\delta\gamma')x''+(\gamma\beta'+\delta\delta')y''
\]
Also enthält $F$ die Form $F''$.

\SourceLeaf{132}
Wegen
\[
(\alpha\alpha'+\beta\gamma')(\gamma\beta'+\delta\delta')
 -(\alpha\beta'+\beta\delta')(\gamma\alpha'+\delta\gamma')
 =(\alpha\delta-\beta\gamma)(\alpha'\delta'-\beta'\gamma')
\]
ist dieser Ausdruck positiv, wenn $\alpha\delta-\beta\gamma$ und $\alpha'\delta'-\beta'\gamma'$ beide positiv oder beide negativ sind; er ist dagegen negativ, wenn eine dieser Zahlen positiv und die andere negativ ist. Die Form $F$ enthält also die Form $F''$ \textit{eigentlich}, wenn $F$ die Form $F'$ auf dieselbe Weise enthält wie $F'$ die Form $F''$, und \textit{uneigentlich}, wenn dies auf verschiedene Weise geschieht.

Daraus folgt: Haben wir beliebig viele Formen $F,F',F'',F'''$ usw., von denen jede die folgende enthält, so enthält die erste die letzte, und zwar \textit{eigentlich}, wenn die Anzahl der Formen, die ihre folgende uneigentlich enthalten, gerade ist, und \textit{uneigentlich}, wenn diese Anzahl ungerade ist.

\textit{Ist die Form $F$ der Form $F'$ und die Form $F'$ der Form $F''$ äquivalent, so ist die Form $F$ der Form $F''$ äquivalent, und zwar eigentlich, wenn die Form $F$ der Form $F'$ auf dieselbe Weise äquivalent ist wie die Form $F'$ der Form $F''$, und uneigentlich, wenn dies auf verschiedene Weise geschieht.}

Da nämlich die Formen $F,F'$ der Reihe nach den Formen $F',F''$ äquivalent sind, enthalten sowohl jene jeweils diese, sodass $F$ die Form $F''$ enthält, als auch diese jene. Also sind $F,F''$ äquivalent. Aus dem Vorigen folgt aber, dass $F$ die Form $F''$ eigentlich oder uneigentlich enthält, je nachdem, ob $F$ der Form $F'$ auf dieselbe oder auf eine andere Weise äquivalent ist wie $F'$ der Form $F''$. Dasselbe gilt dafür, wie $F''$ die Form $F$ enthält. Im ersten Fall sind $F,F''$ also eigentlich, im zweiten uneigentlich äquivalent.

\textit{Die Formen $(a,-b,c)$, $(c,b,a)$, $(c,-b,a)$ sind der Form $(a,b,c)$ äquivalent, und zwar die ersten beiden uneigentlich, die letzte eigentlich.}

Denn $axx+2bxy+cyy$ geht in $ax'x'-2bx'y'+cy'y'$ über, wenn man $x=x'+0.y'$, $y=0.x'-y'$ setzt. Diese Transformation ist wegen $1\times-1-0.0=-1$ uneigentlich. In die Form $cx'x'+2bx'y'+ay'y'$ geht sie durch die uneigentliche Transformation $x=0.x'+y'$, $y=x'+0.y'$ über und in die Form $cx'x'-2bx'y'+ay'y'$ durch die eigentliche $x=0.x'-y'$, $y=x'+0.y'$.

Daraus ist klar, dass jede der Form $(a,b,c)$ äquivalente Form entweder ihr selbst oder der Form $(a,-b,c)$ \textit{eigentlich} äquivalent ist. Enthält eine Form die Form $(a,b,c)$ oder ist sie in ihr enthalten, so enthält sie ebenso entweder die Form $(a,b,c)$ oder die Form $(a,-b,c)$ \textit{eigentlich}, oder sie ist in einer von beiden \textit{eigentlich} enthalten. Die Formen $(a,b,c)$, $(a,-b,c)$ nennen wir \textit{entgegengesetzt}.

\SourceLeaf{133}
\begin{center}
\textit{Benachbarte Formen.}
\end{center}
\Article{160}
Haben die Formen $(a,b,c)$, $(a',b',c')$ dieselbe Determinante und gilt außerdem $c=a'$ und $b\equiv-b'\; (\mathrm{mod.}\,c)$, oder $b+b'\equiv0\; (\mathrm{mod.}\,c)$, so nennen wir diese Formen \textit{benachbart}. Ist eine genauere Bezeichnung nötig, so nennen wir die erste der zweiten \textit{nach links}, die zweite der ersten \textit{nach rechts} benachbart.

So ist \textit{zum Beispiel} die Form $(7,3,2)$ der Form $(3,4,7)$ nach rechts benachbart, die Form $(3.1.3)$ ihrer entgegengesetzten $(3,-1,3)$ nach beiden Seiten.

\textit{Benachbarte Formen sind stets eigentlich äquivalent.} Denn die Form $axx+2bxy+cyy$ geht in die benachbarte Form $cx'x'+2b'x'y'+c'y'y'$ durch die Substitution $x=-y'$, $y=x'+\frac{b+b'}{c}y'$ über, die wegen $0\times\frac{b+b'}{c}-1\times-1=1$ eigentlich ist. Das lässt sich durch Ausmultiplizieren mit Hilfe der Gleichung $bb-ac=b'b'-cc'$ leicht beweisen; $\frac{b+b'}{c}$ ist nach Voraussetzung eine ganze Zahl. --- Diese Definitionen und Folgerungen gelten nicht, wenn $c=a'=0$ ist. Dieser Fall kann aber nur bei Formen eintreten, deren Determinante eine Quadratzahl ist.

Die Formen $(a,b,c)$, $(a',b',c')$ sind eigentlich äquivalent, wenn $a=a'.\ b\equiv b'\; (\mathrm{mod.}\,a)$ gilt. Denn die Form $(a,b,c)$ ist der Form $(c,-b,a)$ eigentlich äquivalent (voriger Artikel), und diese ist der Form $(a',b',c')$ nach links benachbart.

\begin{center}
\textit{Gemeinsame Teiler der Koeffizienten von Formen.}
\end{center}
\Article{161}
\textit{Enthält die Form $(a,b,c)$ die Form $(a',b',c')$, so teilt jeder gemeinsame Teiler der Zahlen $a,b,c$ auch die Zahlen $a',b',c'$, und jeder gemeinsame Teiler der Zahlen $a,2b,c$ auch die Zahlen $a',2b',c'$.}

Geht nämlich die Form $axx+2bxy+cyy$ durch die Substitutionen $x=\alpha x'+\beta y'$, $y=\gamma x'+\delta y'$ in die Form $a'x'x'+2b'x'y'+c'y'y'$ über, so gelten die Gleichungen
\[
\begin{aligned}
a\alpha\alpha+2b\alpha\gamma+c\gamma\gamma&=a'\\
a\alpha\beta+b(\alpha\delta+\beta\gamma)+c\gamma\delta&=b'\\
a\beta\beta+2b\beta\delta+c\delta\delta&=c'
\end{aligned}
\]
aus denen die Aussage sofort folgt. Für ihren zweiten Teil verwendet man anstelle der zweiten Gleichung die Gleichung $2a\alpha\beta+2b(\alpha\delta+\beta\gamma)+2c\gamma\delta=2b'$.

\SourceLeaf{134}
Daraus folgt, dass der größte gemeinsame Teiler der Zahlen $a,b(2b),c$ zugleich den größten gemeinsamen Teiler der Zahlen $a',b'(2b'),c'$ teilt. Enthält also außerdem die Form $(a',b',c')$ die Form $(a,b,c)$, sind die Formen also äquivalent, so ist der größte gemeinsame Teiler der Zahlen $a,b(2b),c$ gleich dem größten gemeinsamen Teiler der Zahlen $a',b'(2b'),c'$, weil dann sowohl jener diesen als auch dieser jenen teilen muss. Haben in diesem Fall die Zahlen $a,b(2b),c$ keinen gemeinsamen Teiler, ist also ihr größter gemeinsamer Teiler $=1$, so haben auch $a',b'(2b'),c'$ keinen gemeinsamen Teiler.

\Needspace{18\baselineskip}
\begin{center}
\textit{Zusammenhang aller gleichartigen Transformationen einer gegebenen Form in eine gegebene Form.}
\end{center}
\Article{162}
\textsc{Aufgabe.} \textit{Enthält die Form}
\[
AXX+2BXY+CYY\quad\mathinner{.\,.\,.}\quad F
\]
\textit{die Form}
\[
axx+2bxy+cyy\quad\mathinner{.\,.\,.\,.\,.}\quad f
\]
\textit{und ist irgendeine Transformation der ersten in die zweite gegeben, so sind daraus alle übrigen ihr gleichartigen Transformationen abzuleiten.}

\textit{Lösung.} Die gegebene Transformation sei $X=\alpha x+\beta y$, $Y=\gamma x+\delta y$. Nehmen wir zunächst an, dass eine andere, ihr gleichartige Transformation $X=\alpha'x+\beta'y$, $Y=\gamma'x+\delta'y$ gegeben ist, um zu untersuchen, was daraus folgt. Bezeichnen wir die Formen und ihre Determinanten durch $F,f,\ =D,d$ und setzen $\alpha\delta-\beta\gamma=e$, $\alpha'\delta'-\beta'\gamma'=e'$, so gilt (Artikel 157) $d=Dee=De'e'$. Da nach Voraussetzung $e,e'$ dasselbe Vorzeichen haben, ist $e=e'$. Es gelten die folgenden sechs Gleichungen:
\[
\begin{aligned}
A\alpha\alpha+2B\alpha\gamma+C\gamma\gamma&=a&\quad\mathinner{.\,.\,.\,.\,.\,.\,.}\quad&[1]\\
A\alpha'\alpha'+2B\alpha'\gamma'+C\gamma'\gamma'&=a&\quad\mathinner{.\,.\,.\,.\,.\,.\,.}\quad&[2]\\
A\alpha\beta+B(\alpha\delta+\beta\gamma)+C\gamma\delta&=b&\quad\mathinner{.\,.\,.\,.\,.\,.}\quad&[3]\\
A\alpha'\beta'+B(\alpha'\delta'+\beta'\gamma')+C\gamma'\delta'&=b&\quad\mathinner{.\,.\,.\,.\,.\,.}\quad&[4]\\
A\beta\beta+2B\beta\delta+C\delta\delta&=c&\quad\mathinner{.\,.\,.\,.\,.\,.\,.}\quad&[5]\\
A\beta'\beta'+2B\beta'\delta'+C\delta'\delta'&=c&\quad\mathinner{.\,.\,.\,.\,.\,.\,.}\quad&[6]
\end{aligned}
\]
Bezeichnen wir der Kürze halber die Zahlen
\[
\begin{gathered}
A\alpha\alpha'+B(\alpha\gamma'+\gamma\alpha')+C\gamma\gamma'\\
A(\alpha\beta'+\beta\alpha')+B(\alpha\delta'+\beta\gamma'+\gamma\beta'+\delta\alpha')+C(\gamma\delta'+\delta\gamma')\\
A\beta\beta'+B(\beta\delta'+\delta\beta')+C\delta\delta'
\end{gathered}
\]
\SourceLeaf{135}
mit $a',2b',c'$, so leiten wir aus den vorigen Gleichungen die folgenden neuen ab\HistoricalNote{*)}{Diese Gleichungen entstehen folgendermaßen: $7$ aus $1.2$, das heißt durch Multiplikation der Gleichung (1) mit der Gleichung (2), oder genauer, indem die linke Seite jener mit der linken Seite dieser und ihre rechte Seite mit der rechten Seite dieser multipliziert und die Produkte gleichgesetzt werden; $8$ aus $1.4+2.3$; die folgende nicht nummerierte aus $1.6+2.5+3.4+3.4$; die folgende nicht nummerierte aus $3.4$; $11$ aus $3.6+4.5$; $12$ aus $5.6$. Eine ähnliche Bezeichnung werden wir auch im Folgenden stets verwenden. Das Ausmultiplizieren müssen wir den Lesern überlassen.}:
\[
\begin{aligned}
a'a'-D(\alpha\gamma'-\gamma\alpha')^2&=aa&\quad\mathinner{.\,.\,.\,.\,.\,.\,.\,.}\quad&[7]\\
2a'b'-D(\alpha\gamma'-\gamma\alpha')(\alpha\delta'+\beta\gamma'-\gamma\beta'-\delta\alpha')&=2ab&\quad\mathinner{.\,.}\quad&[8]\\
4b'b'-D\bigl((\alpha\delta'+\beta\gamma'-\gamma\beta'-\delta\alpha')^2+2ee'\bigr)&=2bb+2ac
\end{aligned}
\]
Addiert man $2Dee'=2d=2bb-2ac$, so ergibt sich
\[
\begin{aligned}
4b'b'-D(\alpha\delta'+\beta\gamma'-\gamma\beta'-\delta\alpha')^2&=4bb&\quad\mathinner{.\,.\,.\,.\,.}\quad&[9]\\
a'c'-D(\alpha\delta'-\gamma\beta')(\beta\gamma'-\delta\alpha')&=bb
\end{aligned}
\]
und durch Subtraktion von $D(\alpha\delta-\beta\gamma)(\alpha'\delta'-\beta'\gamma')=bb-ac$ folgt
\[
\begin{aligned}
a'c'-D(\alpha\gamma'-\gamma\alpha')(\beta\delta'-\delta\beta')&=ac&\quad\mathinner{.\,.\,.\,.\,.\,.}\quad&[10]\\
2b'c'-D(\alpha\delta'+\beta\gamma'-\gamma\beta'-\delta\alpha')(\beta\delta'-\delta\beta')&=2bc&\quad\mathinner{.\,.}\quad&[11]\\
c'c'-D(\beta\delta'-\delta\beta')^2&=cc&\quad\mathinner{.\,.\,.\,.\,.\,.\,.\,.}\quad&[12]
\end{aligned}
\]
Der größte gemeinsame Teiler der Zahlen $a,2b,c$ sei nun $m$, und die Zahlen $\mathfrak{A},\mathfrak{B},\mathfrak{C}$ seien so bestimmt, dass
\[
\mathfrak{A}a+2\mathfrak{B}b+\mathfrak{C}c=m
\]
gilt (Artikel 40). Man multipliziere die Gleichungen $7,8,9,10,11,12$ der Reihe nach mit $\mathfrak{A}\mathfrak{A},\allowbreak 2\mathfrak{A}\mathfrak{B},\allowbreak \mathfrak{B}\mathfrak{B},\allowbreak 2\mathfrak{A}\mathfrak{C},\allowbreak 2\mathfrak{B}\mathfrak{C},\allowbreak \mathfrak{C}\mathfrak{C}$ und addiere die Produkte. Setzen wir nun der Kürze halber
\[
\begin{aligned}
\mathfrak{A}a'+2\mathfrak{B}b'+\mathfrak{C}c'&=T&\quad\mathinner{.\,.\,.\,.\,.\,.\,.\,.}\quad&[13]\\
\mathfrak{A}(\alpha\gamma'-\gamma\alpha')&+\mathfrak{B}(\alpha\delta'+\beta\gamma'-\gamma\beta'-\delta\alpha')&&\\
{}+\mathfrak{C}(\beta\delta'-\delta\beta')&=U&\quad\mathinner{.\,.}\quad&[14]
\end{aligned}
\]
wobei $T,U$ offenbar ganze Zahlen sind, so ergibt sich
\[
TT-DUU=mm
\]
Wir gelangen also zu dem eleganten Schluss, dass \textit{aus je zwei beliebigen gleichartigen Transformationen der Form $F$ in $f$ eine ganzzahlige Lösung der unbestimmten Gleichung $tt-Duu=mm$ folgt}, nämlich $t=T,\ u=U$. Da wir in
\SourceLeaf{136}
unseren Überlegungen nicht vorausgesetzt haben, dass die Transformationen \textit{verschieden} sind, muss sogar \textit{eine} Transformation, zweimal betrachtet, eine Lösung liefern. Dann gilt aber wegen $\alpha'=\alpha,\ \beta'=\beta$ usw. auch $a'=a,\ b'=b,\ c'=c$ und daher $T=m,\ U=0$, eine Lösung, die sich von selbst versteht.

Betrachten wir nun die erste Transformation und die Lösung der unbestimmten Gleichung als bekannt und untersuchen, wie daraus die zweite Transformation abgeleitet werden kann oder wie $\alpha',\beta',\gamma',\delta'$ von $\alpha,\beta,\gamma,\delta,\ T,U$ abhängen. Dazu multiplizieren wir zunächst die Gleichung $[1]$ mit $\delta\alpha'-\beta\gamma'$, $[2]$ mit $\alpha\delta'-\gamma\beta'$, $[3]$ mit $\alpha\gamma'-\gamma\alpha'$, $[4]$ mit $\gamma\alpha'-\alpha\gamma'$ und addieren die Produkte. Das ergibt
\[
(e+e')a'=(\alpha\delta'-\beta\gamma'-\gamma\beta'+\delta\alpha')a\quad\mathinner{.\,.\,.\,.\,.\,.}\quad[15]
\]
Ebenso ergibt sich aus
\[
(\delta\beta'-\beta\delta')([1]-[2])+(\alpha\delta'-\beta\gamma'-\gamma\beta'+\delta\alpha')([3]+[4])+(\alpha\gamma'-\gamma\alpha')([5]-[6])
\]
die Gleichung
\[
2(e+e')b'=2(\alpha\delta'-\beta\gamma'-\gamma\beta'+\delta\alpha')b\quad\mathinner{.\,.\,.\,.\,.}\quad[16]
\]
Schließlich folgt aus $(\delta\beta'-\beta\delta')([3]-[4])+(\alpha\delta'-\gamma\beta')[5]+(\delta\alpha'-\beta\gamma')[6]$:
\[
(e+e')c'=(\alpha\delta'-\beta\gamma'-\gamma\beta'+\delta\alpha')c\quad\mathinner{.\,.\,.\,.\,.\,.}\quad[17]
\]
Setzt man diese Werte $(15,16,17)$ in $13$ ein, so ergibt sich
\[
(e+e')T=(\alpha\delta'-\beta\gamma'-\gamma\beta'+\delta\alpha')(\mathfrak{A}a+2\mathfrak{B}b+\mathfrak{C}c)
\]
oder
\[
2eT=(\alpha\delta'-\beta\gamma'-\gamma\beta'+\delta\alpha')m\quad\mathinner{.\,.\,.\,.\,.\,.\,.}\quad[18]
\]
woraus sich $T$ viel leichter ableiten lässt als aus $[13]$. --- Verbindet man diese Gleichung mit $15,16,17$, so erhält man $ma'=Ta,\ 2mb'=2Tb,\ mc'=Tc$. Setzt man diese Werte von $a',2b',c'$ in die Gleichungen $7\text{---}12$ ein und schreibt für $TT$ den Ausdruck $mm+DUU$, so gehen sie nach den nötigen Änderungen über in
\[
\begin{aligned}
(\alpha\gamma'-\gamma\alpha')^2mm&=aaUU\\
(\alpha\gamma'-\gamma\alpha')(\alpha\delta'+\beta\gamma'-\gamma\beta'-\delta\alpha')mm&=2abUU
\end{aligned}
\]
\SourceLeaf{137}
\[
\begin{aligned}
(\alpha\delta'+\beta\gamma'-\gamma\beta'-\delta\alpha')^2mm&=4bbUU\\
(\alpha\gamma'-\gamma\alpha')(\beta\delta'-\delta\beta')mm&=acUU\\
(\alpha\delta'+\beta\gamma'-\gamma\beta'-\delta\alpha')(\beta\delta'-\delta\beta')mm&=2bcUU\\
(\beta\delta'-\delta\beta')^2mm&=ccUU
\end{aligned}
\]
Mit Hilfe der Gleichung $[14]$ und der Gleichung $\mathfrak{A}a+2\mathfrak{B}b+\mathfrak{C}c=m$ leitet man daraus leicht Folgendes ab: Man multipliziert die erste, zweite und vierte; die zweite, dritte und fünfte; die vierte, fünfte und sechste Gleichung jeweils mit $\mathfrak{A},\mathfrak{B},\mathfrak{C}$ und addiert die Produkte.
\[
\begin{aligned}
(\alpha\gamma'-\gamma\alpha')Umm&=maUU\\
(\alpha\delta'+\beta\gamma'-\gamma\beta'-\delta\alpha')Umm&=2mbUU\\
(\beta\delta'-\delta\beta')Umm&=mcUU
\end{aligned}
\]
Dividiert man dann durch $mU$\HistoricalNote{*)}{Das wäre nicht erlaubt, wenn $U=0$ wäre. Dann aber würde die Richtigkeit der Gleichungen $19,20,21$ sofort aus der ersten, dritten und sechsten der vorangehenden Gleichungen folgen.}, so ergibt sich
\[
\begin{aligned}
aU&=(\alpha\gamma'-\gamma\alpha')m&\quad\mathinner{.\,.\,.\,.\,.\,.\,.\,.\,.}\quad&[19]\\
2bU&=(\alpha\delta'+\beta\gamma'-\gamma\beta'-\delta\alpha')m&\quad\mathinner{.\,.\,.\,.\,.}\quad&[20]\\
cU&=(\beta\delta'-\delta\beta')m&\quad\mathinner{.\,.\,.\,.\,.\,.\,.\,.\,.}\quad&[21]
\end{aligned}
\]
Aus jeder dieser Gleichungen lässt sich $U$ viel leichter als aus $[14]$ ableiten. --- Zugleich folgt daraus, dass bei jeder Wahl von $\mathfrak{A},\mathfrak{B},\mathfrak{C}$, die auf unendlich viele verschiedene Arten möglich ist, sowohl $T$ als auch $U$ jeweils denselben Wert behalten.

Multipliziert man nun die Gleichung $18$ mit $\alpha$, $19$ mit $2\beta$ und $20$ mit $-\alpha$, so ergibt sich durch Addition
\[
2\alpha eT+2(\beta a-\alpha b)U=2(\alpha\delta-\beta\gamma)\alpha'm=2e\alpha'm
\]
Ebenso folgt aus $\beta[18]+\beta[20]-2\alpha[21]$
\[
2\beta eT+2(\beta b-\alpha c)U=2(\alpha\delta-\beta\gamma)\beta'm=2e\beta'm
\]
Ferner folgt aus $\gamma[18]+2\delta[19]-\gamma[20]$
\[
2\gamma eT+2(\delta a-\gamma b)U=2(\alpha\delta-\beta\gamma)\gamma'm=2e\gamma'm
\]
Schließlich ergibt sich aus $\delta[18]+\delta[20]-2\gamma[21]$
\[
2\delta eT+2(\delta b-\gamma c)U=2(\alpha\delta-\beta\gamma)\delta'm=2e\delta'm
\]

\SourceLeaf{138}
Setzt man in diesen Formeln für $a,b,c$ die Werte aus $1,3,5$ ein, so folgt
\[
\begin{aligned}
\alpha'm&=\alpha T-(\alpha B+\gamma C)U\\
\beta'm&=\beta T-(\beta B+\delta C)U\\
\gamma'm&=\gamma T+(\alpha A+\gamma B)U\\
\delta'm&=\delta T+(\beta A+\delta B)U\text{\textsuperscript{*)}}
\end{aligned}
\]
\HistoricalDisplayNote{*)}{Daraus ergibt sich leicht
\[
\begin{aligned}
AeU&=(\delta\gamma'-\gamma\delta')m\\
2BeU&=(\alpha\delta'-\delta\alpha'+\gamma\beta'-\beta\gamma')m\\
CeU&=(\beta\alpha'-\alpha\beta')m
\end{aligned}
\]}

Aus der vorangehenden Untersuchung folgt, dass es keine Transformation der Form $F$ in $f$ gibt, die der gegebenen gleichartig und nicht in der Formel
\[
\begin{aligned}
X&=\tfrac{1}{m}(\alpha t-(\alpha B+\gamma C)u)x+\tfrac{1}{m}(\beta t-(\beta B+\delta C)u)y\\
Y&=\tfrac{1}{m}(\gamma t+(\alpha A+\gamma B)u)x+\tfrac{1}{m}(\delta t+(\beta A+\delta B)u)y\quad\mathinner{.\,.\,.}\quad(\mathrm{I})
\end{aligned}
\]
enthalten wäre. Dabei durchlaufen $t,u$ alle ganzen Zahlen, die der Gleichung $tt-Duu=mm$ genügen. Daraus können wir aber noch nicht schließen, dass alle Werte von $t,u$, die jene Gleichung erfüllen, in die Formel $(\mathrm{I})$ eingesetzt geeignete Transformationen liefern. Es gilt jedoch:

1. Dass die Form $F$ durch eine Substitution aus beliebigen solchen Werten von $t,u$ stets in die Form $f$ übergeht, lässt sich durch Ausmultiplizieren mit Hilfe der Gleichungen $1,3,5$ und der Gleichung $tt-Duu=mm$ leicht bestätigen. Die eher lange als schwierige Rechnung lassen wir der Kürze halber weg.

2. Jede aus der Formel abgeleitete Transformation ist der gegebenen gleichartig. Denn
\[
\begin{aligned}
&\tfrac{1}{m}(\alpha t-(\alpha B+\gamma C)u)\times\tfrac{1}{m}(\delta t+(\beta A+\delta B)u)\\
&\quad-\tfrac{1}{m}(\beta t-(\beta B+\delta C)u)\times\tfrac{1}{m}(\gamma t+(\alpha A+\gamma B)u)\\
&\qquad=\tfrac{1}{mm}(\alpha\delta-\beta\gamma)(tt-Duu)=\alpha\delta-\beta\gamma
\end{aligned}
\]
3. Haben die Formen $F,f$ verschiedene Determinanten, so kann die Formel $(\mathrm{I})$ für manche Werte von $t,u$ Substitutionen liefern, die \textit{Brüche} enthalten und deshalb verworfen werden müssen. Alle übrigen sind jedoch geeignete Transformationen, und weitere gibt es nicht.

4. Haben dagegen die Formen $F,f$ dieselbe Determinante und sind sie daher \textit{äquivalent}, so liefert die Formel $(\mathrm{I})$ keine Transformationen, die Brüche enthalten,
\SourceLeaf{139}
und stellt somit in diesem Fall die vollständige Lösung der Aufgabe dar. Das beweisen wir folgendermaßen.

Aus dem Satz des vorigen Artikels folgt in diesem Fall, dass $m$ zugleich ein gemeinsamer Teiler der Zahlen $A,2B,C$ ist. Wegen $tt-Duu=mm$ gilt $tt-BBuu=mm-ACuu$, also ist $tt-BBuu$ durch $mm$ teilbar. Damit ist erst recht $4tt-4BBuu$ dadurch teilbar, und folglich, weil $2B$ durch $m$ teilbar ist, auch $4tt$ durch $mm$ und daher $2t$ durch $m$. Somit sind $\frac{2}{m}(t+Bu),\ \frac{2}{m}(t-Bu)$ ganze Zahlen, und zwar, weil ihre Differenz $\frac{4}{m}Bu$ gerade ist, entweder beide gerade oder beide ungerade. Wären beide ungerade, so wäre auch ihr Produkt ungerade. Dieses ist aber als Vierfaches der Zahl $\frac{1}{mm}(tt-BBuu)$, die wir gerade als ganzzahlig erkannt haben, notwendig gerade. Dieser Fall ist also unmöglich; daher sind $\frac{2}{m}(t+Bu),\ \frac{2}{m}(t-Bu)$ stets gerade und somit $\frac{1}{m}(t+Bu),\ \frac{1}{m}(t-Bu)$ ganze Zahlen. Daraus folgt ohne Mühe, dass alle vier Koeffizienten in $(\mathrm{I})$ stets ganzzahlig sind. \mbox{\textit{W.z.b.w.}}

Aus dem Vorigen folgt: Kennt man alle Lösungen der Gleichung $tt-Duu=mm$, so lassen sich daraus alle der gegebenen Transformation gleichartigen Transformationen der Form $(A,B,C)$ in $(a,b,c)$ ableiten. Wie man jene Lösungen findet, werden wir im Folgenden zeigen. Hier bemerken wir nur, dass ihre Anzahl stets endlich ist, wenn $D$ negativ oder positiv und zugleich eine Quadratzahl ist, dagegen unendlich, wenn $D$ positiv und keine Quadratzahl ist. Gilt dieser letzte Fall und gilt zugleich die Gleichheit $D$$=d$ nicht (oben Punkt $3^0$), so müsste außerdem untersucht werden, wie sich im Voraus diejenigen Werte von $t,u$, die von Brüchen freie Substitutionen liefern, von denen unterscheiden lassen, die gebrochene ergeben. Für diesen Fall werden wir aber unten eine andere Methode angeben, die diesen Nachteil nicht hat (Artikel 214).

\Needspace{16\baselineskip}
\textit{Beispiel.} Die Form $xx+2yy$ geht durch die eigentliche Substitution $x=2x'+7y'$, $y=x'+5y'$ in die Form $(6,24,99)$ über. Gesucht sind \textit{alle} eigentlichen Transformationen jener Form in diese. Hier ist $D=-2,\ m=3$, also lautet die zu lösende Gleichung $tt+2uu=9$. Ihr wird auf sechs verschiedene Arten genügt, nämlich durch die Wahl $t=3,-3,1,-1,1,-1;\ u=0,0,2,2,-2,-2$ der Reihe nach. Die dritte und die sechste Lösung ergeben Substitutionen mit Brüchen und müssen deshalb verworfen werden. Aus den übrigen folgen vier Substitutionen:
\[
x=\left\{\begin{aligned}
&2x'+7y'\\
&-2x'-7y'\\
&-2x'-9y'\\
&2x'+9y'
\end{aligned}\right.
\qquad
y=\left\{\begin{aligned}
&x'+5y'\\
&-x'-5y'\\
&x'+3y'\\
&-x'-3y'
\end{aligned}\right.
\]
deren erste die gegebene ist.
